% =========================================================
%  sesi/09/sesi.tex
%  Sesi 9 — Integrasi FSM dan A*, Dua Level Berurutan
% =========================================================

\sesi{09}{Integrasi FSM dan A*, Dua Level Berurutan}{2.4}{150 menit}

Minggu 9 adalah pertemuan penutup blok Python. Dua lapisan yang Anda bangun
secara terpisah akhirnya bekerja bersama: mesin keadaan memutuskan apa yang
harus dilakukan penjaga, dan \istilah{pathfinding} menentukan bagaimana
melakukannya. Pemisahan itu penting. Mesin keadaan tidak perlu tahu bentuk
peta, dan pencarian jalur tidak perlu tahu mengapa penjaga ingin bergerak.

Sesi ini juga menyatukan kedua level menjadi satu permainan yang berjalan
tanpa dimulai ulang. Pemain mengumpulkan enam artefak di Level Rak, melewati
pintu, lalu masuk ke Level Lorong tempat penjaga mengejar. Skor dan nyawa
dibawa dari level pertama ke level kedua. Di ujung sesi, tiga peta Lorong
dipakai sebagai tiga \istilah{skenario uji}, dan hasil pengamatannya menjadi
bahan laporan \istilah{playtest} Mini Proyek AI yang dikumpulkan minggu ini.

% ---------------------------------------------------------
\bagian{Tujuan Sesi}

\begin{tujuanbox}
Setelah menuntaskan sesi ini, Anda dapat menunjukkan:
\begin{enumerate}[leftmargin=1.5em, itemsep=0.2em]
  \item satu kali jalan permainan dari Level Rak sampai Level Lorong tanpa
        memulai ulang program;
  \item penjaga yang memanggil A* hanya ketika berada pada keadaan mengejar,
        sedangkan keadaan berdiam dan berpatroli tidak memanggilnya;
  \item \istilah{repath} berkala dengan selang 0,3 detik yang dapat Anda
        tunjukkan pada kode dan hitung pada saat berjalan;
  \item tiga skenario uji yang berjalan dengan pilihan \code{--scenario 1},
        \code{2}, dan \code{3}, beserta catatan waktu dan kesan perilakunya;
  \item perhitungan batas performa penjaga dengan angka 36 piksel, radius
        7 petak, kecepatan patroli 90, dan kecepatan mengejar 165.
\end{enumerate}
\end{tujuanbox}

% ---------------------------------------------------------
\bagian{Prasyarat}

\begin{prasyaratbox}
\begin{itemize}[leftmargin=1.4em, itemsep=0.25em]
  \item Checkpoint sesi 8 berjalan: penjaga menemukan jalan memutar dengan
        A* dan jalurnya tergambar di layar.
  \item Checkpoint sesi 7 berjalan: mesin keadaan penjaga berpindah antara
        \code{idle}, \code{patrol}, dan \code{chase}.
  \item Anda dapat menjalankan program dengan pilihan baris perintah dan
        membaca keluaran \code{argparse}.
  \item Bab XI buku ajar sudah dibaca, khususnya bagian tentang penggabungan
        lapisan, hitung ulang berkala, dan penilaian mutu perilaku.
  \item Kelompok Mini Proyek AI sudah terbentuk, beranggotakan 2--3 mahasiswa.
\end{itemize}
\end{prasyaratbox}

% ---------------------------------------------------------
\bagian{Peta Kemampuan}

\petasesi
  {Pencarian A* berjalan pada satu peta Lorong; Level Rak berakhir sendiri
   dan skor belum dibawa ke level berikutnya.}
  {Keadaan mengejar memanggil A*, \istilah{repath} berkala 0,3 detik,
   peralihan Level Rak ke Level Lorong, dan tiga peta skenario uji.}
  {Mini Proyek AI, Ujian Tengah Semester, dan peralihan ke Unity pada
   pertemuan berikutnya.}

% ---------------------------------------------------------
\bagian{Konsep yang Dipakai}

\subsection*{Keputusan dan rute adalah dua pekerjaan}

\code{GuardBrain} menghasilkan keadaan, bukan gerak. Ketika jarak pemain
masuk radius 7 petak, yaitu 252 piksel, keadaan berubah menjadi
\code{chase}. Perubahan itu hanya menyatakan niat. Pelaksanaannya ada di
\code{Game}: keadaan \code{chase} memanggil \code{_chase}, keadaan
\code{patrol} memanggil \code{_patrol}, dan keadaan \code{idle} tidak
memanggil keduanya.

Pemisahan itu menentukan tempat A* dipakai: hanya pada keadaan mengejar.
Saat berpatroli, penjaga menuju waypoint secara langsung karena rutenya
memang bebas halangan; saat berdiam, tidak ada tujuan yang perlu dicari.
Batas itu menghemat perhitungan dan menjaga perilaku tetap mudah dibaca.

\subsection*{Repath berkala}

A* setiap bingkai itu boros: pada 60 bingkai per detik, satu penjaga
memicu 60 pencarian padahal petak pemain belum tentu berubah. Hasilnya
karena itu disimpan dan dipakai sampai salah satu dari dua keadaan berikut
terjadi.

\begin{enumerate}[leftmargin=1.5em, itemsep=0.2em]
  \item \textbf{Pewaktu habis.} \code{repath_timer} dikurangi
        \code{delta_seconds}; ketika mencapai nol, rute dihitung ulang dan
        pewaktu diisi kembali dengan \code{REPATH_SECONDS = 0.3}. Pada 60
        bingkai per detik, 0,3 detik berarti sekitar 18 bingkai.
  \item \textbf{Jalur kosong.} Jalur yang sudah habis dilalui tidak dapat
        dipakai lagi. Tanpa syarat ini, penjaga berhenti di ujung jalur
        walaupun pemain masih terlihat.
\end{enumerate}

Selama selang itu, rute masih mengarah ke posisi pemain beberapa saat lalu,
sehingga penjaga sempat salah arah. Itulah pertukaran yang harus diterima;
besarnya diatur oleh satu angka.

\subsection*{Menempuh jalur}

Jalur berupa petak, sedangkan penjaga bergerak dalam piksel per detik.
\code{Guard.follow} menuju pusat petak berikutnya dengan vektor
ternormalisasi dikali kecepatan 165. Langkah dibatasi agar tidak melampaui
tujuan; bila jaraknya tidak lebih dari \code{WAYPOINT_REACHED}, yaitu 6
piksel, \code{path_index} dinaikkan.

\subsection*{Batas performa}

Beban pencarian berlipat ketika salah satu faktornya naik: jumlah penjaga,
pencarian per detik, dan petak yang diperiksa sekali pencarian. Dengan satu
penjaga, selang 0,3 detik, dan sekitar 280 petak, bebannya kira-kira 900
petak per detik. Bila dihitung setiap bingkai, angkanya sekitar 16.800.
Selisih itu terasa sebagai laju bingkai yang turun.

\subsection*{Dua level dalam satu program}

Kedua level memakai kamus berbeda, tetapi berbagi skor, nyawa, dan kelas
\code{Game}. Di \code{_update_rak}, pintu yang terbuka dan disentuh pemain
memanggil \code{_load_lorong}. Fungsi itu menyiapkan peta, pemain, penjaga,
dan mesin keadaan baru tanpa mengembalikan skor atau nyawa. Karena masih
dalam satu siklus, jendela tidak pernah tertutup.

\subsection*{Tiga skenario uji}

Peta 1 adalah lorong terbuka dengan dua pasang pilar; jalur lurus hampir
selalu cukup, sehingga hasilnya menjadi pembanding. Peta 2 adalah lorong
berbelok bertingkat yang menuntut jalan memutar. Peta 3 memuat bangunan
besar di tengah, dan di sanalah manfaat A* paling jelas terlihat.

% ---------------------------------------------------------
\bagian{Tutorial}

Kerjakan langkah berikut berurutan. Setiap langkah adalah satu tindakan yang
selesai; jalankan program setiap kali satu bagian baru selesai ditulis.

\begin{langkah}
  \item \textbf{Siapkan folder kerja.} Salin folder proyek sesi 8 menjadi
        folder sesi 9. Jalankan \code{python main.py --level 2} untuk
        memastikan checkpoint lama masih bekerja sebelum ada yang diubah.

  \item \textbf{Tambahkan selang hitung ulang.} Pada \file{settings.py},
        tambahkan \code{REPATH_SECONDS = 0.3} di dekat pengaturan penjaga.
        Pastikan pula \code{CELL = 36}, \code{GUARD_RADIUS_CELLS = 7.0},
        \code{GUARD_PATROL_SPEED = 90.0}, dan
        \code{GUARD_CHASE_SPEED = 165.0} masih bernilai sama.

  \item \textbf{Tambahkan dua peta.} Pada \file{levels.py}, tambahkan
        skenario 2 dan 3 ke dalam kamus \code{SCENARIOS}. Setiap peta tetap
        berukuran 24 kolom dan 14 baris, memuat satu \code{S} dan satu
        \code{E}, serta menggunakan tanda \texttt{\#} untuk dinding.
        Tambahkan pula titik patroli dan titik awal penjaga untuk kedua peta
        itu pada \code{PATROL_POINTS} dan \code{GUARD_START}.

  \item \textbf{Ubah pembangun level.} Ubah \code{build_lorong} agar menerima
        nomor skenario, memilih baris peta yang sesuai, dan mengembalikan
        nama level yang memuat nomor itu. Tanpa nomor, ketiga peta akan
        saling menimpa.

  \item \textbf{Tambahkan pilihan skenario.} Pada \file{main.py}, tambahkan
        argumen \code{--scenario} dengan pilihan 1, 2, dan 3 serta nilai
        bawaan 1. Teruskan nilai itu kepada \code{Game} melalui
        \code{args}.

  \item \textbf{Muat skenario yang diminta.} Pada \code{_load_lorong},
        panggil \code{build_lorong(self.args.scenario)}. Periksa bahwa
        pemain, penjaga, dan titik patroli semuanya berasal dari kamus level
        yang baru, bukan dari peta sebelumnya.

  \item \textbf{Tulis hitung ulang berkala.} Ubah \code{_chase} agar
        mengurangi \code{repath_timer} dengan \code{delta_seconds}, menghitung
        petak pemain, lalu memanggil \code{find_path} hanya bila pewaktu
        sudah habis atau jalur kosong. Sesudah menghitung, isi kembali
        pewaktu dengan \code{S.REPATH_SECONDS}.

  \item \textbf{Sambungkan kedua level.} Pada \code{_update_rak}, ganti
        keadaan menang dengan pemanggilan \code{_load_lorong} ketika pintu
        terbuka dan pemain menyentuhnya. Pastikan skor, nyawa, dan jumlah
        artefak yang tersisa tidak diatur ulang di dalam fungsi itu.

  \item \textbf{Periksa penjaga yang tertangkap.} Ketika pemain tertangkap,
        penjaga harus kembali ke titik awalnya, mesin keadaan dibuat baru,
        dan \code{repath_timer} dinolkan. Tanpa itu, penjaga dapat langsung
        berada pada keadaan mengejar dengan jalur lama.

  \item \textbf{Jalankan alur lengkap.} Jalankan \code{python main.py},
        kumpulkan keenam artefak, lewati pintu, lalu masuk ke Lorong. Mainkan
        sekali lagi dengan \code{python main.py --level 2 --scenario 2} dan
        \code{--scenario 3} untuk memastikan ketiga peta dapat dimuat.
\end{langkah}

\begin{ujibox}
Setelah Langkah 10, periksa lima hal berikut.
\begin{itemize}[leftmargin=1.4em, itemsep=0.2em]
  \item Peralihan terjadi tepat ketika pemain menyentuh pintu yang sudah
        terbuka. Bila pintu terbuka tetapi peralihan tidak terjadi, periksa
        syarat \code{door.open} dan \code{colliderect}.
  \item Skor dan nyawa tidak kembali ke nilai awal setelah pindah level.
        Bila kembali, periksa apakah \code{reset()} ikut dipanggil.
  \item Penjaga hanya memanggil A* ketika mengejar. Bila jalur dihitung saat
        berpatroli, periksa cabang keadaan pada \code{_update_lorong}.
  \item Rute berubah kira-kira setiap 0,3 detik selama pemain bergerak.
        Bila berubah setiap bingkai, periksa nilai \code{REPATH_SECONDS}.
  \item Ketiga skenario dapat dimuat dan masing-masing menampilkan peta yang
        berbeda. Bila peta sama, periksa kunci kamus dan nilai
        \code{args.scenario}.
\end{itemize}
\end{ujibox}

% ---------------------------------------------------------
\bagian{Kode Sesi Ini}

Potongan berikut adalah bagian pokok checkpoint sesi 9. Berkas lengkapnya,
termasuk kelas pemain dan level Rak yang tidak berubah, tersedia pada folder
\file{code/sesi-09}. Bandingkan kode Anda dengan daftar ini setelah program
berjalan, bukan sebelumnya.

\tampilkodebaris[Python]{code/sesi-09/gudang_artefak/levels.py}{67}{99}
  {Peta skenario 2 dan 3 yang dipakai sebagai skenario uji}

\tampilkodebaris[Python]{code/sesi-09/gudang_artefak/levels.py}{101}{134}
  {Titik patroli, titik awal penjaga, dan pembangun level}

\tampilkodebaris[Python]{code/sesi-09/gudang_artefak/game.py}{105}{126}
  {Level Rak dan peralihan ke Level Lorong}

\tampilkodebaris[Python]{code/sesi-09/gudang_artefak/game.py}{165}{178}
  {Keadaan mengejar memanggil A* dan menghitung ulang setiap 0,3 detik}

\checkpoint{09}{Program pada folder itu dapat dijalankan dengan
\code{python main.py}; permainan berpindah dari Level Rak ke Level Lorong
tanpa dimulai ulang, dan ketiga skenario dapat dibuka dengan
\code{--scenario 1}, \code{2}, atau \code{3}.}

% ---------------------------------------------------------
\bagian{Kesalahan Umum}

\begin{galatbox}
\begin{itemize}[leftmargin=1.4em, itemsep=0.35em]
  \item \textbf{Laju bingkai jatuh ketika penjaga mengejar.}
        Gejala: gerak menjadi tersendat dan FPS turun tajam. Penyebabnya,
        A* dijalankan setiap bingkai karena \code{repath_timer} tidak pernah
        diisi kembali atau \code{REPATH_SECONDS} bernilai nol. Periksa baris
        \code{self.repath_timer = S.REPATH_SECONDS} sesudah pencarian.
  \item \textbf{Penjaga berhenti di ujung jalur.}
        Gejala: penjaga diam walaupun pemain masih terlihat. Penyebabnya,
        jalur sudah habis tetapi pencarian ulang belum waktunya. Periksa
        syarat \code{or not self.path} pada \code{_chase}.
  \item \textbf{Peralihan level tidak pernah terjadi.}
        Gejala: pintu terbuka, pemain menyentuhnya, tetapi level tetap.
        Penyebabnya, pemeriksaan pintu berada di cabang yang tidak dieksekusi
        atau \code{_load_lorong} dipanggil sebelum pintu terbuka. Periksa
        urutan pada \code{_update_rak}.
  \item \textbf{Skor kembali nol setelah pindah level.}
        Gejala: HUD menunjukkan skor 0 di Lorong. Penyebabnya,
        \code{_load_lorong} memanggil \code{reset} atau mengatur ulang
        \code{self.score}. Periksa isi fungsi itu.
  \item \textbf{\code{KeyError: 3} ketika memuat skenario.}
        Gejala: program berhenti sebelum jendela menampilkan peta.
        Penyebabnya, kunci 3 belum ada pada \code{SCENARIOS},
        \code{PATROL_POINTS}, atau \code{GUARD_START}. Periksa ketiga kamus
        itu dan ejaan nomornya.
  \item \textbf{Garis jalur tertinggal setelah penjaga berhenti mengejar.}
        Gejala: garis kuning masih tampak padahal keadaan sudah berpatroli.
        Penyebabnya, \code{self.path} tidak dikosongkan ketika keadaan
        berubah. Periksa cabang \code{else} pada \code{_update_lorong} atau
        syarat penggambaran.
\end{itemize}
\end{galatbox}

% ---------------------------------------------------------
\bagian{Lembar Kerja}

Isi tabel pengamatan berikut pada catatan praktikum Anda. Mainkan setiap
skenario dengan urutan pemain yang sama agar hasilnya dapat dibandingkan.
Kolom terakhir memuat kesan Anda, bukan hanya angka.

{\small
\renewcommand{\arraystretch}{1.4}
\begin{tabularx}{\linewidth}{@{}L{2.6cm} C{2.2cm} C{2.0cm} C{1.8cm} Y@{}}
\toprule
\textbf{Skenario} & \textbf{Waktu sampai tertangkap (detik)} &
\textbf{Banyak \istilah{repath}} & \textbf{FPS rata-rata} &
\textbf{Kesan mutu perilaku} \\
\midrule
1. Lorong terbuka & & & & \\[1.0cm]
2. Lorong berbelok & & & & \\[1.0cm]
3. Halangan memutar & & & & \\[1.0cm]
\bottomrule
\end{tabularx}
}

\subsection*{Ketentuan Mini Proyek AI}

Mini Proyek AI dikumpulkan dan dipresentasikan pada minggu 9. Ketentuan
berikut diringkas dari RPS dan Bab XI agar dapat Anda periksa satu per satu.

{\footnotesize
\renewcommand{\arraystretch}{1.25}
\begin{tabularx}{\linewidth}{@{}L{2.7cm} Y@{}}
\toprule
\textbf{Aspek} & \textbf{Ketentuan} \\
\midrule
Kelompok & 2--3 mahasiswa; setiap anggota memahami bagiannya. \\
NPC & FSM yang sekurang-kurangnya memuat \istilah{idle},
\istilah{patrol}, dan \istilah{chase}. \\
Pencarian jalur & A* hanya pada keadaan \code{chase}, dengan peta
berhalangan, \istilah{heuristik}, dan \istilah{antrean prioritas}. \\
Skenario uji & Tiga: lorong terbuka, lorong berbelok, dan halangan yang
memaksa jalur memutar; sertakan angka dan kesan. \\
Laporan \istilah{playtest} & 2--3 halaman; tiga skenario, angka pengukuran,
kesan mutu, perbaikan, dan kendala. \\
Demonstrasi & 10 menit pada minggu 9, dilanjutkan tanya jawab. \\
Keluaran & Kode Python/Pygame, GDD singkat, dan laporan \istilah{playtest}. \\
\bottomrule
\end{tabularx}
}

Penilaiannya memakai bobot berikut: implementasi FSM 25 persen, implementasi
pencarian jalur 25 persen, penggabungan keduanya 20 persen, pengujian beserta
laporan 20 persen, dan presentasi 10 persen. Penggabungan dan pengujian
bersama-sama menyumbang 40 persen, lebih besar daripada komponen kode mana
pun.

% ---------------------------------------------------------
\bagian{Checklist Sesi}

\begin{checklistbox}
\begin{checklist}
  \item Saya dapat menunjukkan peralihan dari Level Rak ke Level Lorong dalam
        satu kali jalan tanpa memulai ulang program.
  \item Saya dapat menjelaskan pekerjaan lapisan keputusan dan lapisan
        navigasi, serta mengapa keduanya dipisahkan.
  \item Saya dapat menunjukkan bahwa A* hanya dipanggil pada keadaan
        mengejar.
  \item Saya dapat menjelaskan dua sebab pencarian diulang: pewaktu 0,3 detik
        habis dan jalur sudah kosong.
  \item Saya dapat menghitung perkiraan beban pencarian per detik untuk satu
        penjaga pada selang 0,3 detik.
  \item Saya telah menjalankan ketiga skenario uji dan mencatat angka serta
        kesan perilakunya.
  \item Saya dapat menyebutkan ketentuan Mini Proyek AI dan memeriksa
        kelengkapan kelompok saya.
\end{checklist}
\end{checklistbox}

% ---------------------------------------------------------
\bagian{Pengait ke Pertemuan Berikutnya}

Blok Python berakhir di sini. Seluruh perkakas yang Anda bangun sejak sesi 1
--- siklus permainan, vektor, fisika, sprite, mesin keadaan, dan pencarian
jalur --- tetap berlaku, tetapi letak tanggung jawabnya akan berpindah.
Minggu 10 adalah Ujian Tengah Semester atas Materi 1--9. Sesudah itu,
pekerjaan berlanjut di Unity: siklus, fisika, tabrakan, dan penyajian
diserahkan kepada mesin, sedangkan Anda memusatkan perhatian pada desain
level, pengalaman bermain, dan penilaian mutu hasilnya.
